\documentclass[11pt,a4paper]{article}
\usepackage[margin=25mm,headheight=15pt]{geometry}
\usepackage[T1]{fontenc}
\usepackage[utf8]{inputenc}
\usepackage{lmodern,amsmath,amssymb,amsthm,mathtools,bm}
\usepackage{xcolor,graphicx,booktabs,tabularx,array,longtable}
\usepackage{microtype,enumitem,fancyhdr}
\usepackage{xurl}
\usepackage{hyperref}
\definecolor{accent}{HTML}{163B55}
\hypersetup{colorlinks=true,linkcolor=accent,citecolor=accent,urlcolor=accent,
 pdftitle={Failure of split-product approximations for time-dependent form sums},pdfauthor={Matthew J. Colbrook},
 pdfsubject={AIM 917: mathematical result and proof}}
\urlstyle{same}
\setlength{\parindent}{0pt}
\setlength{\parskip}{0.35em}
\setlength{\emergencystretch}{3em}
\setlist{itemsep=0.25em,topsep=0.3em}
\newtheorem{theorem}{Theorem}[section]
\newtheorem{lemma}[theorem]{Lemma}
\newtheorem{proposition}[theorem]{Proposition}
\newtheorem{corollary}[theorem]{Corollary}
\theoremstyle{remark}
\newtheorem{remark}[theorem]{Remark}
\newtheorem{definition}[theorem]{Definition}
\DeclareMathOperator{\Tr}{Tr}
\DeclareMathOperator{\supp}{supp}
\DeclareMathOperator{\Lip}{Lip}
\DeclareMathOperator{\Rea}{Re}
\newcommand{\R}{\mathbb R}
\newcommand{\C}{\mathbb C}
\newcommand{\E}{\mathbb E}
\newcommand{\N}{\mathbb N}
\newcommand{\ip}[2]{\langle #1,#2\rangle}
\pagestyle{fancy}
\fancyhf{}
\fancyhead[L]{\small\sffamily AIM 917}
\fancyhead[R]{\small\sffamily Matthew J. Colbrook}
\fancyfoot[C]{\thepage}
\renewcommand{\headrulewidth}{0.3pt}
\allowdisplaybreaks[2]
\title{\textbf{Failure of split-product approximations for time-dependent form sums}}
\author{Matthew J. Colbrook\\[3pt]
{\small Department of Applied Mathematics and Theoretical Physics}\\
{\small University of Cambridge, United Kingdom}\\
{\small\href{mailto:mjc249@maths.cam.ac.uk}{\texttt{mjc249@maths.cam.ac.uk}}}}
\date{3 November 2026}
\begin{document}
\maketitle
\begin{abstract}
We give counterexamples to the stated exponential and resolvent split-product formulas for nonautonomous Schrodinger evolution. Narrow time spikes cause failure of the proposed limits while retaining the required common form domain and time regularity. A separate exponential counterexample uses uniformly bounded perturbations.
\end{abstract}
\noindent\textbf{Repository target.}
This manuscript concerns AIM problem 917 at revision
\texttt{5d7e26e8444e}. The precise scope is stated below.
\medskip

\section{Main result and proof}
\label{sec:917}

\textbf{Outcome.} The two assertions in Problem~917 are false under its
stated hypotheses. The counterexample below has a bounded, finite-rank,
self-adjoint perturbation at every individual time. The restricted forms
are continuously differentiable in operator norm from the common form
domain to its dual. Thus no convention concerning non-semibounded closed
forms is involved. These assumptions allow time dependence that is
smooth on the form scale while remaining badly behaved on fixed
Hilbert-space vectors. The final strengthening even has uniformly
bounded perturbations. The precise question is in the catalogue~\cite{AIM}; its
source is \cite[conditions (V), (V$'$), and Remark~(3), equations
(32)--(33)]{VV2011}.

\begin{theorem}\label{thm:917}
There exist a positive self-adjoint operator $H_0$ on a separable complex
Hilbert space $\mathcal H$ and a family of bounded finite-rank
self-adjoint operators $V(t)$, $0\le t\le1$, satisfying every hypothesis of
Problem~917, for which the exponential split products at $(t,s)=(1,0)$
have no strong limit. The resolvent split products do not converge
strongly, or weakly, to the specified unitary propagator $U(1,0)$.
\end{theorem}

\begin{proof}
We first specify the sampling scales. Set
\[
 \alpha=\frac\pi2,\qquad \delta=\frac1{64},\qquad
 \beta(s)=\begin{cases}(1-s^2)^2,&|s|\le1,\\0,&|s|>1.\end{cases}
\]
The function $\beta$ is $C^1$, satisfies $\beta(0)=1$ and
$0\le\beta\le1$, and has $|\beta'|\le2$. Choose positive integers $N_n$
recursively, beginning with $N_1\ge4$, with $N_n\ge4N_{n-1}$. Once $N_n$
is chosen, define
\[
 h_n=N_n^{-1},\qquad \tau_n=1-h_n,\qquad g_n=\alpha N_n,
 \qquad w_n=\frac{2^{-n}}{1024g_n},
 \qquad v_n(t)=g_n\beta\!\left(\frac{t-\tau_n}{w_n}\right).
\]
In choosing $N_n$, require also
\begin{equation}\label{eq:917-riemann}
 \frac1{N_n}\sum_{j=0}^{N_n-1}\sum_{k<n}v_k(j/N_n)\le\delta.
\end{equation}
This is a valid recursive choice: the finite sum $\sum_{k<n}v_k$ is
already determined and continuous, and
\[
 \int_0^1\sum_{k<n}v_k(t)\,dt
 \le\sum_{k<n}2g_kw_k
 \le\frac1{512}<\delta.
\]
Its left Riemann sums converge to this integral. Consequently arbitrarily
large integers meet \eqref{eq:917-riemann} and the growth requirement.

Write $I_n=(\tau_n-w_n,\tau_n+w_n)$. Since $w_n<h_n/10$ and
$h_{n+1}\le h_n/4$, these intervals are disjoint. In fact all $I_k$ with
$k>n$ lie strictly to the right of $\tau_n$. The uniform $N_n$-grid used
in the question therefore samples $I_n$ exactly once, at its last point
$\tau_n=(N_n-1)/N_n$, and samples none of the later intervals. The endpoint
$1$ is not sampled.

Let $\mathcal H=\ell^2(\mathbb N_0)$ with orthonormal basis
$e_0,e_1,\ldots$. Define a diagonal positive self-adjoint operator by
\[
 H_0e_0=e_0,\qquad H_0e_n=\lambda_ne_n,
 \qquad \lambda_n=\left(\frac{g_n}{w_n^2}\right)^2\quad(n\ge1).
\]
Thus $\lambda_n\ge1$, and its form domain and form are
\[
 Q=D(H_0^{1/2}),\qquad
 q_0[u]=|u_0|^2+\sum_{n\ge1}\lambda_n|u_n|^2.
\]
Let $S_n=|e_0\rangle\langle e_n|+|e_n\rangle\langle e_0|$, zero on the
orthogonal complement, and put
\[
 V(t)=\begin{cases}v_n(t)S_n,&t\in I_n,\\0,&t\notin\bigcup_n I_n.
 \end{cases}
\]
In particular $V(1)=0$. At every time, $V(t)$ is a bounded finite-rank
self-adjoint operator. Its closed Hermitian form has domain all of
$\mathcal H$.

We check the uniform form assumptions explicitly. For $u\in Q$ and
$t\in I_n$,
\begin{align}
 |q_{V(t)}[u]|
 &\le\frac{g_n}{\sqrt{\lambda_n}}
       \beta\!\left(\frac{t-\tau_n}{w_n}\right)
       \bigl(|u_0|^2+\lambda_n|u_n|^2\bigr)\notag\\
 &\le w_n^2q_0[u].\label{eq:917-formbound}
\end{align}
Use $q_0^{1/2}$ as the norm on $Q$. Under the natural identifications of
$Q$ and its dual, the restricted form operator is represented by
\[
 B(t)=H_0^{-1/2}V(t)H_0^{-1/2}
     =w_n^2\beta\!\left(\frac{t-\tau_n}{w_n}\right)S_n
       \quad(t\in I_n),
\]
and by zero outside the intervals. Hence
\begin{equation}\label{eq:917-derivative}
 \|B'(t)\|\le2w_n\qquad(t\in I_n).
\end{equation}
Both $\beta$ and $\beta'$ vanish at the interval endpoints, so $B$ is
norm $C^1$ away from $1$. At the accumulation point,
\[
 \frac{\|B(t)\|}{1-t}
 \le\frac{w_n^2}{h_n-w_n}
 \le\frac{2w_n^2}{h_n}\longrightarrow0
 \quad(t\in I_n,\ n\to\infty),
\]
while \eqref{eq:917-derivative} tends to zero. It follows that $B$ is
norm $C^1$ on $[0,1]$, with $B'(1)=0$. Equivalently,
$t\mapsto\widetilde V(t)$ is norm $C^1$ as a map into
$\mathcal L(Q,Q^*)$. Every $w_n$ is smaller than $1/8$, so
\eqref{eq:917-formbound}--\eqref{eq:917-derivative} imply both of the
catalogue's bounds with $a=1/4$ and $b=0$. Moreover
\[
 q_t[u]=q_0[u]+q_{V(t)}[u]
 \ge\frac34q_0[u]\ge\frac34\|u\|^2.
\]
The standing positivity condition holds with $c=3/4$. The hypotheses
therefore provide the unitary form propagator $U(t,s)$.

We next examine the exponential product on the grid of size $N_n$.
Denote by $K_n^E$ the product of its first $N_n-1$ factors, omitting the
last factor at $\tau_n$. All perturbations sampled in $K_n^E$ come from
the earlier spikes $k<n$. Telescoping the unitary factors and using
$\|e^{-ihV}-I\|\le h\|V\|$ gives
\begin{equation}\label{eq:917-prefixE}
 \|K_n^E-e^{-i\tau_nH_0}\|
 \le h_n\sum_{j=0}^{N_n-2}\|V(jh_n)\|\le\delta.
\end{equation}
At the last grid point, $V(\tau_n)=g_nS_n$ and $h_ng_n=\pi/2$, whence
\[
 e^{-ih_nV(\tau_n)}e_0=-ie_n.
\]
Since $H_0e_0=e_0$, \eqref{eq:917-prefixE} implies
\[
 P_{N_n}^Ee_0=c_ne_n+r_n,
 \qquad c_n=-ie^{-i\tau_n}e^{-ih_n\lambda_n},
 \qquad |c_n|=1,\quad\|r_n\|\le\delta.
\]
Consequently, whenever $m\ne n$,
\[
 \|P_{N_n}^Ee_0-P_{N_m}^Ee_0\|\ge\sqrt2-2\delta>0.
\]
The exponential products are not strongly Cauchy, and therefore have no
strong limit.

Finally, let $R_n=(I+ih_nH_0)^{-1}$ and let $K_n^R$ be the analogous
resolvent prefix with the last factor omitted. All factors are
contractions, and
$\|(I+ihV)^{-1}-I\|\le h\|V\|$. The same telescoping argument yields
\[
 \|K_n^R-R_n^{N_n-1}\|\le\delta.
\]
Thus $K_n^Re_0=b_ne_0+s_n$ with $|b_n|\le1$ and $\|s_n\|\le\delta$.
On the two-dimensional span of $e_0,e_n$,
\[
 (I+i\alpha S_n)^{-1}e_0
   =\frac{e_0-i\alpha e_n}{1+\alpha^2},
 \qquad
 \|(I+i\alpha S_n)^{-1}e_0\|=(1+\alpha^2)^{-1/2}.
\]
The final $R_n$ is a contraction, so
\[
 \|P_{N_n}^Re_0\|\le(1+\alpha^2)^{-1/2}+\delta<1.
\]
As $U(1,0)$ is unitary,
\[
 \|P_{N_n}^Re_0-U(1,0)e_0\|
 \ge1-(1+\alpha^2)^{-1/2}-\delta>0.
\]
This excludes the proposed strong limit. It excludes weak convergence to
$U(1,0)$ as well: for contractions $P_n$ and a unitary $U$, weak
convergence $P_n\to U$ implies strong convergence, because
\[
 \|P_nx-Ux\|^2
 \le2\|x\|^2-2\operatorname{Re}\langle P_nx,Ux\rangle\longrightarrow0.
\]
\end{proof}

\paragraph{An independent construction of the propagator.}
In this example the existence of a unitary propagator can also be seen
without an abstract form-evolution theorem. Indeed,
\[
 \int_0^1\|V(t)\|\,dt
 =\left(\int_{-1}^1(1-s^2)^2\,ds\right)\sum_{n\ge1}g_nw_n
 =\frac{16}{15}\frac1{1024}=\frac1{960}.
\]
The operator-valued function is Bochner measurable: on each spike it is
a continuous scalar multiple of a fixed finite-rank operator, and its
range lies in a separable subspace of $\mathcal B(\mathcal H)$.
The interaction-picture perturbation
$W(t)=e^{itH_0}V(t)e^{-itH_0}$ is also Bochner measurable and has the
same integrable norm. On $I_n$ its two matrix entries merely acquire
the continuous phases $e^{\pm it(\lambda_n-1)}$.
The norm-convergent Dyson series therefore solves
\[
 Z(t,s)=I-i\int_s^t W(\tau)Z(\tau,s)\,d\tau.
\]
It is unitary because $W(t)$ is self-adjoint almost everywhere. Hence
$e^{-itH_0}Z(t,s)e^{isH_0}$ is a unitary propagator. It agrees with the
form propagator: testing the form equation against finite-support
vectors transported by $e^{-itH_0}$ cancels the free generator and gives
the same integral equation; its solution is unique. By density this
identifies the propagators on all of $\mathcal H$. Unitarity in the
integral equation gives the stronger quantitative statement
\[
 \|U(1,0)-e^{-iH_0}\|=\|Z(1,0)-I\|\le\frac1{960}.
\]
Thus the actual evolution is very close to the free evolution even
though the selected grids produce the failures proved above.

The obstruction is an exact sampling effect. Each narrow time spike is
small, together with its
first derivative, on the form scale, but its peak acts strongly on a
fixed low-energy vector in the Hilbert space. A selected uniform grid
hits that peak and transfers the vector to a new high-energy basis
vector. Earlier sampled spikes have uniformly small total effect by
construction. The example even has a common operator domain
$D(H_0)$ for all $H_0+V(t)$, since each $V(t)$ is bounded. The form
hypotheses do not control the action of the sampled perturbations on
fixed Hilbert-space vectors. The following strengthening shows that a
uniform operator bound alone still does not ensure the exponential
limit.

\subsection{The exponential product can fail even for uniformly bounded perturbations}
\label{sec:917bounded}

The exponential obstruction has a stronger version. A uniform bound on
$\|V(t)\|$ alone does not repair the statement; time regularity in the
Hilbert-space operator topology matters as well.

\begin{theorem}\label{thm:917bounded}
For every $\varepsilon>0$, there is a family satisfying all hypotheses of
Problem~917, with every $V(t)$ self-adjoint and finite rank, such that
\[
 \sup_{0\le t\le1}\|V(t)\|\le1,
 \qquad \int_0^1\|V(t)\|\,dt\le\varepsilon,
\]
but its exponential split products at $(t,s)=(1,0)$ have no strong limit.
The restricted forms are continuously differentiable in operator norm,
and $V(t)$ is continuous in the weak operator topology.
\end{theorem}

\begin{proof}
Fix $\delta=1/100$ and put
$\varepsilon_0=\min\{\varepsilon,\delta/8\}$.
Let $\beta(s)=(1-s^2)^2$ on $[-1,1]$ and zero elsewhere, as before.
We recursively choose integers $N_n$, finite sets
$J_n\subset\{0,\ldots,N_n-1\}$, and widths $w_n>0$.

Suppose the earlier groups have been chosen. Let $E_{n-1}$ be the
union of their closed bump intervals, and let
\[
 Z_{n-1}=\{0,1\}\cup
  \bigcup_{k<n}\{j/N_k:0\le j\le N_k-1\}.
\]
The set $E_{n-1}$ is a finite union of intervals; its length will be at
most $\varepsilon_0$. Its indicator is Riemann integrable. Choose $N_n$
sufficiently large that
\[
 \frac1{N_n}\#\{0\le j<N_n:j/N_n\in E_{n-1}\}\le\frac\delta4,
 \qquad
 \frac{\# Z_{n-1}}{N_n}\le\frac\delta4.
\]
This is possible because $|E_{n-1}|\le\varepsilon_0\le\delta/8$ and
$Z_{n-1}$ is finite. Require also $N_n>N_{n-1}$ when $n>1$, and take
$1/N_n$ smaller than every positive point of
$E_{n-1}\cup Z_{n-1}$. The latter requirement is possible because this
finite union of closed intervals and points stays a positive distance
from zero after zero itself is removed. Set
\[
 J_n=\{0\le j<N_n:j/N_n\notin E_{n-1}\cup Z_{n-1}\},
 \qquad m_n=\#J_n.
\]
Then $m_n/N_n\ge1-\delta/2$, and $1\in J_n$. Choose
$w_n<1/(4N_n)$ so small that the closed
intervals
\[
 [j/N_n-w_n,j/N_n+w_n],\qquad j\in J_n,
\]
are pairwise disjoint, lie in $(0,1)$, and avoid
$E_{n-1}\cup Z_{n-1}$. Also impose
\[
 2N_nw_n\le\varepsilon_0\,2^{-n}.
\]
All these requirements are possible because there are finitely many
chosen centers, each at positive distance from the closed forbidden set.
They ensure that all bump intervals in all groups are disjoint and that
\[
 |E_{n-1}|\le\sum_{k<n}2m_kw_k\le\varepsilon_0,
\]
justifying the induction. Define
\[
 b_n(t)=\sum_{j\in J_n}\beta\!\left(\frac{t-j/N_n}{w_n}\right).
\]
Each $b_n$ is $C^1$, with $0\le b_n\le1$ and
$\|b_n'\|_\infty\le2/w_n$. Every later group vanishes at every point
of the $N_n$-grid, by its avoidance of $Z_{k-1}$.

On $\ell^2(\mathbb N_0)$ let $H_0e_0=e_0$ and
$H_0e_n=\lambda_ne_n$, choosing
\[
 \lambda_n=1+2\pi k_nN_n,\qquad k_n\in\mathbb N,
\]
so large that
\[
 \lambda_n^{-1/2}(1+2/w_n)\le2^{-n}/16.
\]
Let $S_n=|e_0\rangle\langle e_n|+|e_n\rangle\langle e_0|$, and put
\[
 V(t)=\sum_{n\ge1}b_n(t)S_n.
\]
At each time at most one summand is nonzero. Thus $V(t)$ is bounded,
self-adjoint and of rank at most two, with $\|V(t)\|\le1$. Moreover,
\[
 \int_0^1\|V(t)\|\,dt
 \le\sum_{n\ge1}2m_nw_n\le\varepsilon_0\le\varepsilon.
\]
For $Q=D(H_0^{1/2})$, the restricted form is represented by
\[
 B(t)=H_0^{-1/2}V(t)H_0^{-1/2}
     =\sum_{n\ge1}\lambda_n^{-1/2}b_n(t)S_n.
\]
Both this series and its derivative series converge uniformly in
operator norm, and
$\sup_t\|B(t)\|\le1/16$, $\sup_t\|B'(t)\|\le1/16$.
Hence $B$ is norm $C^1$. The form bounds hold, for example, with
$a=1/8$ and $b=0$, while
$q_0+q_{V(t)}\ge(15/16)q_0\ge(15/16)\|\cdot\|^2$.
This verifies every form hypothesis. Matrix elements are continuous
first for vectors in $Q$, by form continuity, and then for arbitrary
Hilbert-space vectors by density and the uniform bound on $\|V(t)\|$.
Thus $V$ is weak-operator continuous. It is not strongly continuous
at zero: $V(0)e_0=0$, while $V(1/N_n)e_0=e_n$ and $1/N_n\to0$.
Its Bochner measurability and the existence of the unitary propagator
follow as in the preceding Dyson construction.

Consider the product on the $N_n$-grid, with $h_n=1/N_n$. At a node
$j\in J_n$, its perturbation is exactly $S_n$. At all other nodes it is
an earlier perturbation or zero; no later group contributes. Compare the
product with the ideal one that keeps $S_n$ at the nodes of $J_n$ and
uses zero at all other nodes. Telescoping products of unitaries and using
$\|e^{-ihV}-I\|\le h\|V\|$ bounds their operator-norm difference by
\[
 \frac1{N_n}\#\{0\le j<N_n:j/N_n\in E_{n-1}\}
 \le\frac\delta4\le\delta.
\]
On $\operatorname{span}\{e_0,e_n\}$, the free step is precisely
$e^{-ih_n}I$, since
$h_n(\lambda_n-1)=2\pi k_n$. It therefore commutes with each retained
$S_n$-step on this subspace. Writing $a_n=m_n/N_n\in[1-\delta,1]$,
the ideal product sends $e_0$ to
\[
 v_n=e^{-i}\bigl(\cos(a_n)e_0-i\sin(a_n)e_n\bigr).
\]
Consequently $\|P_{N_n}^Ee_0-v_n\|\le\delta$. For $m\ne n$,
\[
 \|P_{N_n}^Ee_0-P_{N_m}^Ee_0\|
 \ge\sqrt2\sin(1-\delta)-2\delta>0.
\]
The sequence is not strongly Cauchy. Finally, the Dyson estimate gives
$\|U(1,0)-e^{-iH_0}\|\le\varepsilon_0$, so the actual evolution can be
made arbitrarily close to the free evolution while the splitting still
has this fixed nonconvergence gap.
\end{proof}



\begin{thebibliography}{99}
\small
\setlength{\itemsep}{3pt}

\bibitem{AIM}
MColbrook/AIM, \emph{AIM 917: Split products for time-dependent form-sum Schrödinger evolution},
repository statement at revision \texttt{5d7e26e8444e}.
\url{https://github.com/MColbrook/AIM/blob/5d7e26e8444eb98de91988454564e45e011107c4/problems/917-split-products-for-time-dependent-form-sum-schrodinger-evolution.md}.

\bibitem{VV2011}
P.-A.~Vuillermot and W.~F.~Wreszinski,
\emph{Product approximations for solutions to a class of evolution
equations in Hilbert space},
Portugaliae Math. \textbf{68} (2011), 317--343.
\href{https://doi.org/10.4171/PM/1894}{doi:10.4171/PM/1894}.
\end{thebibliography}
\end{document}
